% Page 027

\seite{27}

\jahresueberschrift{Das Schuljahr 1899}
\pstart
Aufgenommen wurden sieben Neulinge, sodaß\ob
die Schule im kommenden Schuljahr von 50\ob
Kindern, 21 Mädchen und 29 Knaben,\ob
besucht wird. Ein Schulkind wurde der\ob
Schule zu Pronsfeld, Kreis Prüm, über\ohb
wiesen.

\pend
\abschnitt{Vertretung}
\pstart

\margmark{Erlaß 11.9.99          D\\I Fleuth}%
a der Lehrer des Nachbarortes Seffern zu einer\ob
zehnwöchentlichen Übung einberufen ist, ist dessen\ob
Vertretung dem Lehrer der hiesigen Schule übertragen.\ob
Aus diesem Grunde ist in der hiesigen Schule Halb\ohb
tagsunterricht eingeführt mit der Maßnahme, daß ab\ohb
wechselnd der Unterricht bald Vormittag, bald Nach\ohb
mittag in der hiesigen Schule abgehalten wird. –\ob
Da der Lehrer in Seffern (Kohlbrecher) nach 14\ob
Tagen für dienstuntauglich erklärt wurde,\ob
konnte er seinen Dienst in Seffern wieder\ob
aufnehmen und somit hörte die Vertretung\ob
auf.

Ausfall des Unterrichts wegen Hitze.

Im Monat August u.\ September mußte der\ob
Unterricht wiederholt wegen zu großer Hitze\ob
ausgesetzt werden.
\pend
